
\documentclass[10pt, a4paper]{amsart}


\usepackage{amsmath,amssymb,amscd,amsfonts}

\newtheorem{thm}{Theorem}[section]
\newtheorem{lem}[thm]{Lemma}
\newtheorem{rem}[thm]{Remark}
\newtheorem{prop}[thm]{Proposition}
\newtheorem{cor}[thm]{Corollary}
\newtheorem{defn}[thm]{Definition}
\newtheorem{assu}[thm]{Assumption}
\newtheorem{claim}[thm]{Claim}
\newtheorem{ex}[thm]{Example}
\newtheorem{pb}[thm]{Problem}

\newcommand\Pic{{\text{\rm Pic}}}
\newcommand\mult{{\text{\rm mult}}}
\newcommand\alb{{ a}}
\newcommand\lra{\longrightarrow}
\newcommand\Alb{{ A}}
\newcommand\Mod{\text{\rm Mod}}
\newcommand\ot{{\otimes}}
\newcommand\OO{{\mathcal{O}}}
\newcommand\PP{{\mathcal{P}}}
\newcommand\QQ{{\mathcal{Q}}}
\newcommand\FF{{\mathcal{F}}}
\newcommand\LL{{\mathcal{L}}}
\newcommand\PPP{{\mathbb{P}}}
\newcommand\Z{{\mathbb{Z}}}
\newcommand\Q{{\mathbb{Q}}}
\newcommand\CC{{\mathbb{C}}}
\newcommand\C{{\mathbb{C}}}
\newcommand\II{{\mathcal{I}}}
\newcommand\sh{{\hat {\mathcal S}}}
\newcommand\s{{ {\mathcal S}}}
%%%%%%%%%%%%%%%%%%%%%%%%%
%%%%%%%%%%%%%%%%%%%%%%%%%
\title{A conjecture of Ueno}
\author {Christopher D. Hacon}
\date{}
\begin{document}
\abstract We prove the following conjecture of Ueno: {\it  
If $X$ is a smooth complex projective variety with
$\kappa (X)= 0$, and $F$ is the general fiber of the Albanese morphism
$\alb :X\lra \Alb $, then $\kappa (F)=0$.} As a consequence of this, we show 
that if $f:X\lra Y$ is an algebraic fiber space with $Y$ of maximal
Albanese dimension, then $\kappa (X)\geq \kappa (F)+\kappa (Y)$.
\endabstract

\maketitle
%{\bf PRELIMINARY VERSION. PLEASE DO NOT CIRCULATE! Questions, comments and suggestions welcome!}
%%%%%%%%%%%%%%%%%%%%%%%%%%%%%%%%%%%
\section{Introduction}
%%%%%%%%%%%%%%%%%%%%%%%%%%%%%%%%%%%
In this paper we prove part of a long standing conjecture of Ueno
which is an important test case of the well known conjecture $C_{n,m}^+$ of 
Iitaka (cf. \cite{Mo} \S 7).

\bigskip
\noindent {\bf Conjecture K} (Ueno).
{\em Let $X$ be a smooth complex projective variety with $\kappa
(X)=0$, and let $\alb :X\lra \Alb $ be its Albanese morphism. Then
\begin{enumerate}
\item $\alb$ is surjective and has connected fibers, i.e. $\alb$ is an 
algebraic fiber space;
\item if $F$ is the general fiber of $\alb$, then $\kappa (F)=0$;
\item there is an \'etale covering $B\lra \Alb$ such that 
$X\times _\Alb B$ is birationally equivalent to $F\times B$ over $B$.
\end{enumerate} }

In \cite{Ka1}, Kawamata proved $(1)$ above.
In this paper, we will prove $(2)$. 
As a consequence of this, we will show that Iitaka's Conjecture
$C_{n,m}$ holds when the base has maximal Albanese dimension.
Since any curve of genus $g\geq 1$ has maximal Albanese dimension,
in particular this recovers $C_{n,1}$ which was proven by Kawamata
in \cite{Ka2}.
We remark that a weaker but related statement was obtained in \cite{CH1}
Theorem 4, 
via the use of generic vanishing theorems.
In this paper, we rely on the generalization of the generic vanishing
theorems of Green and Lazarsfeld obtained in \cite{Hac} which implies a 
non-vanishing result for coherent sheaves on an abelian variety.
In order to apply this non-vanishing result to higher direct images
of powers of dualizing sheaves, we will make use of
certain asymptotic multiplier ideal sheaves
which are a variant of the ones introduced in \cite{Hac}.
We will also make use of a result of Viehweg regarding weak positivity of 
push-forwards of powers of dualizing sheaves, and of a generalization of
a result of Simpson (cf. \cite{CH2} Theorem 3.2) 
regarding the cohomological support
loci of powers of dualizing sheaves. 
\bigskip


\noindent{\bf Acknowledgments.} The author would like to thank J. A. Chen for 
useful discussions. The author was partially supported
by NSA research grant no: MDA904-03-1-0101 and by a grant from the
Sloan Foundation.


\medskip
\noindent{\bf Notation and conventions.}
We work over the field of complex numbers.
We identify Cartier divisors and line bundles on
a smooth variety, and we use the additive and
multiplicative notation interchangeably. If
$X$ is a smooth projective variety, we let $K_X$ be a  canonical divisor,
so that $\omega _X=\OO_X(K_X)$, and we denote by
$\kappa(X)$ the Kodaira dimension,  by
$q(X):=h^1(\OO_X)$ the {\em irregularity} and by $P_m(X):=h^0(\omega _X^{\ot m
})$
the {\em $m-$th plurigenus}. 
We denote by  $\alb \colon X\to \Alb$ the Albanese map  and by 
$\hat{A}$ the dual abelian variety to $\Alb $ which parameterizes all
topologically trivial line bundles on $X$.
For a $\Q -$divisor $D$
we let $\lfloor D\rfloor$ be the integral part and
$\{D\}$ the fractional part. Numerical
equivalence is denoted by $\equiv$ and we write $D\prec E$ if
$E-D$ is an effective divisor. If $f\colon X\to Y$ is a
morphism of normal varieties, we write $K_{X/Y}:=K_X-f^*K_Y$ and $F_{X/Y}$ denotes
the generic fiber of $f$. 
The rest of the notation is
standard in algebraic geometry.
%%%%%%%%%%%%%%%%%%%%%%%%%%%%%%%%%%%
\section{Preliminaries }
%%%%%%%%%%%%%%%%%%%%%%%%%%%%%%%%%%%%%
Here we recall several results from the literature and we prove some technical
statements that will be needed later.
%%%%%%%%%%%%%%%%%%%%%%%%%%%%%%%%%%%%%%%%%%%
\subsection{Non-vanishing for coherent sheaves on an abelian variety}
%%%%%%%%%%%%%%%%%%%%%%%%%%%%%%%%%%%%%%%%%%%
Let $A$ be an abelian variety and $\hat{A}=\Pic ^0(A)$ its 
dual abelian variety. We will denote by $p_A,p_{\hat {A}}$ the projections
of $A\times \hat{A}$ on to $A,\hat {A}$. Let $\PP $ be the normalized 
Poincar\'e line bundle on $A\times \hat {A}$.
Let $L$ be an ample line bundle on $\hat {A}$, then we define
$\hat{L}:=p_{A,*}(p_{\hat{A}}^*L\ot \PP )$. If $\phi _L :\hat{A}\lra A$
denotes the isogeny defined by $\phi _L(\hat{a})=t_{\hat{a}}^*L^\vee\ot L$,
then $\phi _L^*\hat{L}^\vee=L^{\oplus {h^0(L)}}$.
We will need the following result:
\begin{prop}\label{gvt} 
Let $\FF$ be a non-zero
coherent sheaf on an abelian variety $A$ such that for any ample
line bundle $L$ on $\hat {A}$ one has 
$$h^i(\hat {A},\phi _L^*\FF \ot \hat{L}^\vee)=0\ \ \ \forall \ i>0,$$
then $h^0(\FF \ot P)\ne 0$ for some $P\in \Pic ^0(A)$.
\end{prop}
\begin{proof} By Corollary 3.2 (1) of \cite{Hac}, if $h^0(\FF \ot P)=0$ for all
$P\in \Pic ^0(A)$, then $h^i(\FF \ot P)=0$ for all $i\geq 0$ and all
$P\in \Pic ^0(A)$. By a result of Mukai (cf. \cite{Mu}), this implies that
$\FF=0$. 
\end{proof}
We remark that since $\hat{L}^\vee$ pulled back via the finite \'etale
morphism $\phi _L$, splits as a direct sum of ample line bundles,
for the purposes of vanishing theorems (cf. Theorem \ref{omegavan}), 
it behaves as an ample line bundle.

%%%%%%%%%%%%%%%%%%%%%%%%%%%%%%%%%%%%%%%%%%%
\subsection{Multiplier ideal sheaves}\label{MIS}
%%%%%%%%%%%%%%%%%%%%%%%%%%%%%%%%%%%%%%%%%%%
%%%%%%%%%%%%%%%%%%%%%%%%%%%%%%%%%%%%%%%%%%%
%\subsection{Definitions}
%%%%%%%%%%%%%%%%%%%%%%%%%%%%%%%%%%%%%%%%%%%
In this section we recall some well known properties of multiplier
ideal sheaves. We refer to \cite{La} for a more complete treatment.

Let $X$ be a smooth projective variety and $|L|$ a (non-empty)
linear series on $X$, let $\nu : \tilde {X}\lra X$
be log resolution of $|L|$ i.e. a proper birational morphism from a 
smooth projective
variety such that $\nu ^*|L|=|M|+F$ where $|M|$ is base point free and
$F\cup \{ Exceptional \ locous\ of (\nu )\}$ has simple normal crossings support.
For any rational number $c>0$, define the multiplier ideal sheaf
associated to $c$ and $|L|$ by
$$\II (X,c \cdot |L|):= \nu _* \OO _{\tilde{X}}
(K_{\tilde {X}/X}-\lfloor c F \rfloor ).$$
When $D=|L|$ this is just the usual multiplier ideal sheaf associated
to the effective $\Q$-divisor $c\cdot D$, usually denoted by  $\II (c\cdot D)$.
Roughly speaking, the ideals sheaves 
$\II (X,c \cdot |L|)$ are a measure of the singularities
of a general divisor of $|L|$. The more singular divisors corresponding to
deeper ideals. Their definition does not depend  on the choice of a log 
resolution. 
Recall the following well known properties (cf. \cite{La} Example 9.3.5,
Proposition 9.2.32 and Example 9.2.29):
\begin{prop} \label{ideal}
Let $Z$ be an irreducible subvariety of $X$ of codimension $c$ and $D$ a $\Q$ divisor
with $\mult _ZD\geq c$ (i.e. vanishing at a generic point of $Z$ to order 
at least $c$), then 
$\II (D)\subset \II _Z$.
\end{prop}
\begin{prop} \label{ideal2}
Let $|A|,|V|,|W|$ be linear series such that $|A|$ is base point free.
If $|V|\subset |W|$ then $\II (|V|)\subset \II (|W|)$.
If $|W|+|A|\subset |V|$, then
for any positive rational number $c>0$, one has $\II (c\cdot |W|)
\subset \II (c\cdot |V|)$.
\end{prop}
A {\it graded linear series} corresponding to a divisor $L$
is a sequence $V_{\bullet}=\{ V_k \}$
where $k>0$ and each $V_k$ is a finite dimensional subspace of 
$H^0(X,\OO _X(kL))$ such that the natural maps induce inclusions
$$V_m\cdot V_l \subset V_{l+m}\ \ \ for\ all\ l,m>0 .$$
For any rational number $c>0$, the {\it asymptotic multiplier ideal sheaf} 
of such a graded linear series $\II (X,c\cdot V_{\bullet})=\II (X,c\cdot
||V_1||)$ 
is the unique maximal element among the multiplier ideals
$\II (X,\frac {c}{p}|V_{p}| )$. 
The above definition only makes sense if $V_k\ne 0$ for some (and hence
for all sufficiently divisible) integer $k>0$.
%\begin{ex} Let $L$ be a line bundle on $X$ and $Z$ an irreducible
%subvariety of $X$. Let
%\end{ex} 
%%%%%%%%%%%%%%%%%%%%%%%%%%%%%%%%%%%%%%%%%%%
\subsection{Direct images of dualizing sheaves}\label{dualizing}
%%%%%%%%%%%%%%%%%%%%%%%%%%%%%%%%%%%%%%%%%%%
\begin{thm} \label{omegavan}Let $f\colon X\to  Y$ be a  surjective 
map of projective
varieties, $X$ smooth. Let $M$ be a line bundle on $X$
such that $M-\Delta \equiv f^* L $, where $L$
is a $\Q-$divisor on $Y$ and $\Delta $ is an effective $\Q$-divisor on $X$. 
Then
\begin{itemize}
\item[a)] $R^jf_*(\omega _X\ot M\ot \II (\Delta )) $ 
is torsion free for $j\geq 0$;

\item[b)] If $L$ is nef and big, then $H^i(Y,R^jf_*(\omega _X\ot M
\ot \II (\Delta )))
=0$ for $i>0$, $j\geq 0$;

\item[c)] If $L$ is nef and big and $H$ is very ample on $Y$, 
then $$R^jf_*(\omega _X \ot M\ot \II (\Delta ))\ot H^{\ot l}$$ is generated
by global sections for any integer $l\geq \dim Y +1$.
\end{itemize}
\end{thm}
\begin{proof} This result is certainly well known. We briefly indicate
how to recover the above version from the references available in the 
literature. When $Y$ is normal and $(X,\Delta)$ is klt (i.e.
$\II (\Delta )=\OO _X$), statements a) and  b)
are \cite{Ko1} 
Corollary 10.15. Following the proof of \cite{HP} Lemma 2.7,
one sees that the assumption that $Y$ is normal can be dropped.
For the general case, let $\nu :\tilde{X}\lra X$ be a log resolution
of $(X,\Delta)$ so that $\nu ^* \Delta$ has simple normal crossings support.
We notice that $K_ {\tilde{X}}+\nu ^* M -\lfloor \nu ^* \Delta
\rfloor \equiv K_ {\tilde{X}}+(f \circ \nu)^*L+\{\nu ^* \Delta \}$,
and $(\tilde{X}, \{\nu ^* \Delta \})$ is klt.
Therefore, one has that $$\nu _*(\omega _{\tilde{X}}\ot \nu ^* M(-\lfloor \nu ^* \Delta
\rfloor ))=\omega _X \ot M \ot \II (\Delta )$$
and all higher direct images are torsion free and hence vanish.
So, $$R^jf_*(\omega _X\ot M\ot \II (\Delta ))=
R^j(f \circ \nu)_*(\omega _{\tilde{X}}\ot \nu ^* M(-\lfloor \nu ^* \Delta
\rfloor )).$$
Statements a) and b) now follow.
Statement c) follows easily by restricting to general hyperplane sections in
$|H|$.
\end{proof}
%%%%%%%%%%%%%%%%%%%%%%%%%%
\section{Main result}
%%%%%%%%%%%%%%%%%%%%%%%%%%%%%%
\begin{thm} \label{T1}
Let $X$ be a smooth complex projective variety, $a:X\lra A$ a surjective
morphism with connected fibers to an abelian variety. If $\kappa (X)=0$
then $\kappa (F)=0$ where $F$ denotes a generic fiber of $a$.
\end{thm}
\begin{proof}
Let $n:=\dim X$, and $q:=\dim A$.
Assume that $P_1(X)=1$ and $P_N(F)\geq 2$.
For a fixed very ample divisor $H$ on $A$, any $D\in |NK_F|$ and
any rational number $\epsilon>0$, we claim that the homomorphism
$$
\phi _{l,\epsilon H}:
H^0(X,\OO _X(l(NK_X+a^*\epsilon H)))\lra H^0(\OO _A/\II _F^2\ot \OO _A
(kNK_X+a^*\epsilon H))$$ is also surjective for some $l> 0$.
In fact,
by a result of Viehweg (c.f. \cite{V1}, \cite{V2}, see also \cite{Ko2} \S 10),
the homomorphism
$$
H^0(X,\OO _X(l(NK_X+a^*\frac{\epsilon}{2} H)))\lra H^0(F,\OO _F
(kNK_F))$$ is surjective for some $l> 0$.
If $F=F_a$, following \cite{PP}, one sees that the morphism
$$H^0(A, a_* (\omega _X^{lN}))\ot H^{\ot l\epsilon})\lra H^0(a_* (\omega _X^{lN})
\ot \OO _A /m_a^2)$$
is surjective and the claim follows.

Let $$|V _{l,\epsilon H}^{D  }|:=\phi _{l,\epsilon H}^{-1}(lD)$$
where by abuse of notation, 
$\phi _{l , \epsilon H}$ also denotes the corresponding map of linear series. 
Notice that for all $j>0$, there is a natural morphism
$|V^{D}_{l, \epsilon H}|^{\times j}\lra |V^{D}_{lj,\epsilon H}|$.
By the above claim, for all $l$, we have
$V _{l,\epsilon H}^{D}\ne 0$ and the general divisor in $|V _{l,\epsilon H}^{D}|$
intersects $F$ transversely.

We consider the graded linear series
$V^{D}_{\bullet, \epsilon H}:=\{ V^{D}_{l,\epsilon H} \}$
and, for any integer $c>0$, the corresponding asymptotic multiplier ideal sheaf
$$\II ^{c\cdot D}_{\epsilon H}:=\II (X,c\cdot V^{D}_{\bullet, \epsilon H})
.$$
By Section \ref{MIS}, $\II ^{c\cdot {D}}_{\epsilon H}=
\II (\frac{c}{l}|V^{D}_{l,\epsilon H}|) $
for all $l\gg 0$ sufficiently divisible.
We now let $c=q+1$. Sections of $V^{D}_{l,\epsilon H}$ intersect $F$ transversely
and so they have multiplicity at least
$l$ along $D$. Since the codimension of $D$ in $X$ is $q+1$,
by Proposition \ref{ideal},
one has that $$\II ^{c\cdot {D}}_{\epsilon H}\subset \II _{D}.$$
Let $$\FF _\epsilon :=a_*(\omega _X ^{\ot 1+cN}\ot 
\II ^{c\cdot {D}}_{\epsilon H}).$$
We claim that there exists a rational number $1<\epsilon <\frac {1}{c}$
such that for all rational numbers $0<\epsilon '<\epsilon$, one has
$$\FF _\epsilon =\FF _{\epsilon '}.$$
To see this, observe that for all $l\gg 0$ sufficiently divisible, 
there is a homomorphism of
linear series:
$$|V^{{D}}_{l,\epsilon ' H}|\times a^*|l(\epsilon -\epsilon ')H|\lra
|V^{D}_{l,\epsilon  H}|$$
and since $a^*|l(\epsilon -\epsilon ')H|$ is base point free, 
by Proposition \ref{ideal2}, one has
$$\FF _\epsilon\subset\FF _{\epsilon '}
\ \ \ for\ all\ 0<\epsilon '<\epsilon .$$ 
Since $\FF _\epsilon:= a_*(\omega _X \ot M \ot \II (\Delta ))$
where $M-\Delta\equiv{ -c\epsilon }a^*H$ 
and $\Delta\equiv c(NK_X+\epsilon a^*H)$ with 
$c\epsilon<1$, 
by Theorem \ref{omegavan}, 
one sees that  the sheaf 
$$\FF _\epsilon\ot
H^{\ot q+2}$$ 
is generated by global sections (and similarly when $\epsilon$ is replaced by
a smaller rational number $\epsilon '$).
If the inclusion $\FF _{\epsilon '}\subset \FF _\epsilon$
is strict, the number of global sections must 
decrease. It follows that there are at most finitely many 
strict inclusions and so all of these sheaves coincide for all
rational numbers $0<\epsilon '< \epsilon \ll 1$.

We claim that $\FF _{\epsilon}$ satisfies the hypothesis of
Proposition \ref{gvt}. To this end, fix any ample line bundle $L$ on
$\hat{A}$ and pick $0<\epsilon ' \ll 1$ so that $\phi _L^*H-\epsilon ' H$
is ample. By Theorem \ref{omegavan} one sees that $$H^i(\FF _{\epsilon }\ot
\hat {L}^\vee)= H^i(\FF _{\epsilon '}\ot
\hat {L}^\vee)=0$$
as required.

By Proposition \ref{gvt}, one has that 
$$h^0(\omega _X^{\ot 1+ckN}\ot \II _{D_0}\ot a^*P)\geq
h^0(a_*(\omega _X^{\ot 1+ckN}\ot \II ^{c\cdot D_0}_{\epsilon H})\ot P)
=h^0(\FF _\epsilon \ot P)>0$$
for some $P\in \Pic ^0(A)$. We may assume that $a^*P=\OO _X$.
In fact, by Theorem 3.2 of \cite{CH2}, for any integer $m>0$,
the irreducible components of
$$V^0(mK_X):=\{ P\in \Pic ^0(A) | h^0(\omega ^{\ot m}_X \ot a^* P)\ne 0 \}$$
are torsion translates of abelian subvarieties of $\Pic ^0(X)$.
Therefore, we may 
assume that if $a^*P\ne \OO _X$, then
$a^*P\in \Pic ^0(X)_{tors}-\{\OO _X \}$. But since $\kappa (X)=0$ and $P_1(X)>0$,
this is easily seen to be impossible.
Since $\kappa (X)=0$, we have that $h^0(\omega _X^{\ot 1+cN})=1$ 
and so $D\subset Bs |(1+cN)K_X|$. 
Since $P_N(F)\geq 2$, then the divisors $D\in |NK_F|$ cover
$F$ and so $F\subset Bs|(1+cN)K_X|$ but this is impossible as $F$ is a general
fiber.

We now consider the case in which $P_1(X)=0$.
By a well known result of Fujita,
there exists a generically finite cover 
$\tilde{X}\lra X$ such that $\kappa (\tilde{X})
=0$ and $P_1(\tilde {X})=1$.
Let $\tilde{X}\lra \tilde {A}\lra {A}$ be the Stein factorization.
Then, by Kawamata's result (i.e. the first part of Conjecture K), 
$\tilde{X}\lra \tilde {A}$
is surjective with connected fibers and $\tilde {A}$ is birational to
an abelian variety. Let $\tilde{F}$ be the corresponding general fiber.
Since there is a generically finite morphism $\tilde{F}\lra F$, one has that
$0=\kappa (\tilde {F})\geq \kappa (F)$. As $\kappa (X)\geq 0$, one sees that
$\kappa (F)\geq 0$ and the assertion follows.
\end{proof}
%\begin{lem} Let $f:X\lra Y$ be an algebraic fiber space with generic fiber 
%$F_{X/Y}$. 
%If $Y$ is of maximal Albanese dimension and
%$\kappa (X)= \kappa (Y)$, then $\kappa (F_{X/Y})=0$.
%\end{lem}
%\begin{proof} Let $X\lra V$ and $Y\lra W$ be the Iitaka fibrations of $X,Y$.
%By the proof of Theorem 6 of \cite{CH1}, we may assume that (after
%replacing $X,Y,V,W$ by appropriate birational models), the morphism
%$X\lra W$ factors through a birational morphism $V\lra W$. 
%$F_{X/V}=F_{X/W}\lra F_{Y/W}$ is also an algebraic fiber space with generic
%fiber $F_{X/Y}$. Since $F_{Y/W}$ is birational to an abelian variety and
%$\kappa (F_{X/V})=0$, one sees that $\kappa (F_{X/Y})=0$.
%\end{proof} 
\begin{cor} Let $f:X\lra Y$ be an algebraic fiber space with $Y$ of maximal
Albanese dimension, then
$$\kappa(X)\geq \kappa (Y)+\kappa (F_{X/Y}).$$
\end{cor}
\begin{proof} This follows from an argument of Kawamata 
(cf. \cite{Mo} Theorem 6.4).
For the convenience of the reader, we illustrate this, following
the exposition of \cite{Fu} Lemma 4.2  (to which we refer the reader 
interested in further details).
Let $\alpha :Y\lra A$ be the Albanese morphism of $Y$, 
$\varphi :X\lra V$ be the Iitaka fibration of $X$. By definition, this means 
that $\kappa(X)=\dim V$ and
$\kappa (F_{X/V})=0$. Let $$f':F_{X/V}\lra B:=f(F_{X/V})$$ be the induced map.
As $B$ is of maximal Albanese dimension, by Kawamata's Theorem
it follows easily that $B$ is
biratonal to an abelian variety which is an \'etale cover
of an abelian subvariety $B'$ of $A$. 
%Since abelian subvarieties are rigid, 
%we may assume $B'$ is fixed.
By Theorem \ref{T1}, the fibers of $f'$
which we denote by $F'$, satisfy $\kappa (F' )=0$. By a lemma of Kawamata
(cf. \cite{Fu} Lemma 4.3), there exists a subvariety $\Gamma$ of $V$
and a projective variety $\tilde {W}$ and morphisms
$\Phi :\tilde{W }\lra \Gamma$ and $\Psi:\tilde {W}\lra Y$ such that $\Psi$ is
generically finite and for general $v\in \Gamma$, we have that
$\Phi ^{-1}(v)=f(\varphi ^{-1}(v))$ is birational to $B$.
One has that
$$\kappa (Y)\leq \kappa (\tilde {W})\leq \kappa (\Phi ^{-1}(v))+
\dim \Gamma 
\leq \kappa (B)+\dim Y-\dim B=\dim Y-\dim B.$$
Since $F'=f^{-1}(y)\cap \varphi ^{-1}(v)$, 
one may think of $F'$ as the general fiber
of $X_y\lra \varphi (X_y)$.
By the easy addition formula:
$$\kappa (X_y)\leq \kappa (F')+\dim (\varphi (X_y))=\dim (\varphi (X_y)).$$
One sees easily that $$\dim (\varphi (X_y))=\dim X_y-\dim F' =
\kappa (X)+\dim B-\dim Y.$$ 
Therefore we have that
$$\kappa (X) \geq \kappa(X_y)+\dim Y-\dim B\geq \kappa (Y) +\kappa (X_y).$$
\end{proof}
\begin{thebibliography}{Ka1}

\bibitem[CH1]{CH1}
{\sc Chen, J.\,A.} and {\sc Hacon, C.\,D.}:
On algebraic fiber spaces over varieties of maximal albanese dimension,
Duke Math. J., Vol. 111, No. 1, 2002

\bibitem[CH2]{CH2}
{\sc Chen, J.\,A.} and {\sc Hacon, C.\,D.}:
On the irregularity of the Iitaka fibration,
Comm. Algebra {\bf 32} (2004) no. 1, 203-215.

\bibitem[Fu]{Fu}
{\sc Fujino, O.}:
Algebraic fiber spaces whose general  fibers
are of maximal albanese dimension,
Nagoya Math. J.  {\bf 172}  (2003), 111--127.

\bibitem[Hac]{Hac}
{\sc Hacon, C.\,D.}:
A derived category approach to generic vanishing theorem,
to appear in J. Reine Angew. Math 

\bibitem[HP]{HP}
{\sc Hacon, C.\,D.} and {Pardini, R.}:
On the birational geometry of varieties of maximal Albanese dimension.  
J. Reine Angew. Math.  {\bf 546}  (2002), 177--199.

\bibitem[Ka1]{Ka1}
{\sc Kawamata, Y.}:
Characterization of abelian varieties,
Compositio Math. {\bf 43}, 1981, 253-276.

\bibitem[Ka2]{Ka2}
{\sc Kawamata, Y.}: Kodaira dimension of algebraic fiber spaces over a curve,
Invent. Math. {\bf 66} (1982) no. 1, 57-71.


\bibitem[Ko1]{Ko1}
{\sc Koll\'ar, J.}:
Shafarevich maps and automorphic forms,
M. B. Porter Lectures, Princeton Univ. Press, Princeton, 1995

\bibitem[Ko2]{Ko2}
{\sc Koll\'ar, J.}:
Shafarevic maps and plurigenera of algebraic varieties, Invent. Math {\bf 113},
1993, 177-215

\bibitem[La]{La}
{\sc Lazarsfeld, R.}:
Positivity in Algebraic Geometry, II. To appear in 
Ergenbisse der Mathematik (Springer).

\bibitem[Mo]{Mo}
{\sc Mori, S.}:
Classification of higher dimensional varieties, Algebraic Geometry 
(Brunswick, Maine, 1985), Proc. Sympos. Pure Math. {\bf 46},
Amer. Math. Soc., Providence, 1987, 269-331

\bibitem[Mu]{Mu}
{\sc Mukai, S.}:
Duality between $D(X)$ and $D(\hat {X})$
with its application to Picard sheaves,
Nagoya math. J. {\bf 81},
1981, 153-175

\bibitem[PP]{PP}
{\sc}

\bibitem[V1]{V1}
{\sc Viehweg, E.}:
Weak positivity and the additivity of the Kodaira dimension
of certain fiber spaces.
Adv. Studies in Pure Math {\bf 1}, 1983, 329-353

\bibitem[V2]{V2}
{\sc Viehweg, E.}:
Quasi-projective moduli for polarized manifolds,
Ergenbisse der Mathematik {\bf 30} (Springer, 1995)
\end{thebibliography}
\bigskip
\bigskip

\begin{minipage}{13cm}
\parbox[t]{12.5cm}{Christopher D. Hacon\\
Department of Mathematics\\
University of Utah\\
155 South 1400 East, JWB 233\\
Salt Lake City, Utah 84112-0090, USA\\
hacon@math.utah.edu}
\end{minipage}

\end{document}
\end

